\section{Introduction}\label{sec:intro}

\textit{Pleurotus ostreatus} is among the most widely cultivated edible fungi and a workhorse
model for white-rot lignocellulose degradation. Its enzymatic repertoire---laccases,
manganese and versatile peroxidases, lytic polysaccharide monooxygenases and an extensive
glycoside hydrolase complement---underpins both its ecological role and a growing set of
biotechnological applications. Interest has recently extended to controlled-environment and
bioregenerative life support contexts, where fungal biomass offers a protein source that can
be produced from lignocellulosic residues in a closed system. The mycoponic culture format
used here, in which mycelium is grown on micro-structured ceramic tubes supplied with liquid
nutrient medium through an antimicrobial size-exclusion interface, was developed for exactly
this purpose \cite{porterfield2026mycoponics}.

Understanding how such a system works requires resolving what different parts of the mycelium
are doing. Filamentous fungal colonies are not homogeneous: they differentiate into
morphologically and physiologically distinct regions, and in this culture system they form
several visually distinguishable tissue types, including structures that produce and exude
liquid droplets.

For multi-condition designs of this kind, 3$'$-end tag counting has become an attractive
option: it sequences a single fragment per transcript, so cost per sample falls sharply and
library complexity requirements are lower than for full-length protocols. The trade-offs are
less widely appreciated than the benefits. Because one read corresponds to one molecule
regardless of transcript length, no length normalisation applies and TPM or FPKM values are
meaningless. Because reads pile up at the polyadenylation site, quantification depends on the
accuracy of annotated 3$'$ ends rather than on gene bodies. And because the method primes on
poly(A) tracts, it is sensitive to internal priming and to any failure of poly(A) selection.

Those dependencies interact badly with a second issue: the state of public fungal genome
annotation. Genome annotations are produced to describe protein-coding gene models. Ribosomal
DNA is frequently omitted, being repetitive, hard to assemble and of little interest for gene
finding; organellar rRNA is omitted more often still. For a genome browser this is
inconsequential. For read quantification it is not, because a feature that does not exist
cannot receive counts and cannot be filtered out. Current consensus RNA-seq pipelines,
including the NASA GeneLab pipeline used as the basis for this work
\cite{genelab7101G}, implement rRNA removal by dropping features whose annotated biotype is
rRNA---a strategy that is silently ineffective against rRNA that was never annotated.

Both problems are present in an acute form in the dataset analysed here, which allows us to
quantify costs that are usually invisible and to provide the corrected resources that the
species has lacked.

\section{Results}

\subsection{Reference selection changes mapping by an order of magnitude}\label{sec:refs}

Four \textit{P.\ ostreatus} assemblies were compared empirically rather than assumed
(Table~\ref{tab:refs}, Fig.~\ref{fig:refs}a). Uniquely mapped fraction differed far more than
assembly statistics would suggest. The RefSeq-designated reference genome for the species
(PC9) performed \emph{worst}, at 5.45\% uniquely mapped against 24.94\% for the best
candidate---less than a quarter. BOM\_ss5 and BOM\_ss14 differed by 0.02 percentage points,
consistent with their being the two nuclei of a single dikaryon, matching the cultivar used
here.

On matched 3$'$-extended annotations, BOM\_ss5 assigned 6.9\% more reads than PC9.15 and did
so for every one of the 16 libraries, with gains from 4.0\% to 13.5\%
(Fig.~\ref{fig:refs}b). BOM\_ss5 was used for all subsequent analysis; the complete parallel
PC9.15 analysis is retained in the repository.

\subsection{The libraries are dominated by rRNA from two distinct sources}\label{sec:rrna}

Against barrnap-derived loci, nuclear rDNA accounted for a mean of 53.5\% of reads
(range 19.9--69.7\%; Fig.~\ref{fig:fate}b). This is a direct measurement; a $k$-mer probe
estimate made earlier in the analysis (12--46\%) was a substantial underestimate, and we note
this because probe-based estimates are common in quality-control reporting.

Clustering intergenic reads from pooled alignments revealed the second source. In BOM\_ss5,
40.2\% of pooled aligned reads fall in a tandem block structure at
\texttt{JBQVBD010000012.1:2{,}298{,}112--2{,}346{,}320} (the nuclear rDNA array) and 31.3\% on
the mitochondrial contig \texttt{CM148777.1}---approximately 71\% of aligned reads in
identifiable rDNA loci (Supplementary Fig.~S7). barrnap's fragmentary 5S/5.8S/18S/28S calls
fall inside every block, confirming their identity. In PC9.15 the equivalent mitochondrial
block contained 60.9\% of all pooled aligned reads and was confirmed by BLAST at 100\%
identity to the \textit{P.\ ostreatus} mitochondrion.

\textbf{Neither the nuclear rDNA array nor the mitochondrial rRNA is annotated as a feature in
any of the four public annotations.} In the GTF, the mitochondrial rRNA region is flanked only
by tRNA gene models (Fig.~\ref{fig:refs}c). The practical consequence is severe: these reads
cannot be counted, cannot be removed by biotype-based rRNA filtering, and are silently
reported as ``no feature''.

\subsection{Assembly choice changes rDNA behaviour, not just contiguity}

BOM\_ss5 resolves multiple tandem copies of the nuclear rDNA repeat; PC9.15 collapses them.
The consequence is quantitative, not cosmetic: in the collapsed assembly, rDNA reads
distribute across fewer loci and are reported as multi-mapping and discarded, whereas in the
resolved assembly they pile into an identifiable block that can be annotated and accounted
for. This is the mechanism underlying the assignment-rate difference in Section~\ref{sec:refs},
and it means that assembly selection for RNA-seq should consider rDNA representation, which is
not reported in standard assembly metrics.

\subsection{Extending gene 3$'$ ends does not fix the problem}\label{sec:utr}

A natural response to low assignment in 3$'$-tag data is that annotated 3$'$ UTRs are too
short. This is superficially well supported here: of the four annotations, only PC9.15 carries
meaningful UTRs (median 3$'$ UTR 78~bp; the other three have 3~bp, i.e.\ the stop codon only).

We extended gene 3$'$ ends strand-aware by up to 500~bp with a neighbour cap (91.3\% of 13,556
genes extended; median extension 323~bp; 3.79~Mb added). The effect on gene assignment was
15.5\%~$\rightarrow$~15.6\% (Supplementary Fig.~S2).

The diagnostic that distinguishes the two causes is simple and we recommend it as routine:
\emph{cluster the reads that fail to be assigned and ask whether they are diffuse or focal}.
UTR truncation produces diffuse loss distributed just downstream of many genes; unannotated
rRNA produces a small number of very high-coverage blocks. Here a single locus held 60.9\% of
aligned reads. Adding rRNA features reduced ``no feature'' from 30--78\% to 0--1.2\%.

\subsection{Effective yield and the cost of retaining noise-dominated libraries}

Despite 92,430,743 raw reads, only 2,475,808 were assigned to non-rRNA features across all 16
libraries; 7,601,974 were assigned to rRNA features. Effective mRNA yield ranged from 0.09\%
to 10.49\% of raw reads, and only 3 of 16 libraries exceeded 100,000 assigned non-rRNA counts
(Table~\ref{tab:budget}, Fig.~\ref{fig:fate}).

Retaining the lowest-yield libraries actively degraded the analysis. With all 16, the first
principal component correlated with sequencing depth at $r = 0.686$---the dominant axis of
variation was library size, not biology. Excluding the four libraries below 15,000 assigned
non-rRNA counts reduced this to $r = -0.003$, while increasing genes passing expression
filtering from 1,199 to 1,666 and genes with a significant tissue effect from 82 to 281
(Fig.~\ref{fig:pca}). All four tissues retained $n \geq 2$.

\subsection{Resources generated}

\textbf{Annotation.} rRNA-complete, 3$'$-aware GTFs for BOM\_ss5 and PC9.15, with rRNA blocks
declared in per-reference tables.

\textbf{Functional annotation.} Of 12,521 BOM\_ss5 proteins, 5,163 (41\%) have a Swiss-Prot
hit; 39,607 GO/KEGG term assignments; 423 proteins carrying a CAZy-class domain (GH 233,
AA 86, GT 57, PL 30, CBM 13, CE 4); 576 proteins predicted secreted (Supplementary Fig.~S4).
The auxiliary-activity complement---laccases, peroxidases, lytic polysaccharide
monooxygenases---is consistent with the oak sawdust and microcrystalline cellulose in the
growth medium.

\textbf{Metabolic reconstruction.} The first draft genome-scale metabolic model for
\textit{P.\ ostreatus}: 1,915 reactions, 2,028 metabolites and 881 genes, with
gene--protein--reaction rules traceable to supporting proteins. Constraining exchanges to the
mycoponic medium and adding transport reactions raised the number of reactions able to carry
flux from 288 to 2,262. Targeted gapfilling then completed the cofactor pathways with three
added reactions carrying no gene evidence; biotin remained unreachable and was supplied by the
medium rather than given a biosynthesis route, which is biologically appropriate for a likely
auxotroph. The final model has 5,247 reactions of which 2,287 carry flux, all six tested
cofactors producible, and non-zero biomass flux on the defined medium
(Table~\ref{tab:s7}, Supplementary Fig.~S5).

\subsection{The pipeline recovers coherent signal where power permits}

We assessed tissue signal by two independent criteria.

\textbf{Replicate-supported markers.} Tissue specificity was scored with the $\tau$ index on
linear normalised counts. Because $\tau$ is computed on tissue means, a single outlying
replicate can manufacture an apparent marker; we therefore required that \emph{every}
replicate of a tissue exceed the maximum value in any sample of any other tissue. Of the top
50 markers per tissue, this was satisfied by 37 for exudophore ($n=2$), 22 for nodule ($n=3$),
1 for exuding mycelium ($n=3$) and 1 for fuzzy mycelium ($n=4$) (Supplementary Fig.~S3). The
test becomes more stringent as $n$ increases, so these counts are not comparable across
tissues; the qualitative separation is nonetheless clear.

\textbf{Co-expression modules.} WGCNA (12 libraries, 4,748 genes) produced 41 modules, 37 with
at least one enriched GO or KEGG term. Of 168 module\,$\times$\,tissue correlations, 17 reached
$p<0.05$ (8.4 expected by chance) and six survived Benjamini--Hochberg correction---all
associated with exudophore or nodule (Table~\ref{tab:modules}, Fig.~\ref{fig:wgcna}a).

The two methods agree completely on which tissues carry signal, despite being sensitive to
different failure modes. Neither supports exuding or fuzzy mycelium at the achieved depth.

\textbf{Convergent signatures.} The nodule-associated \texttt{brown} module enriches for
extracellular space (7/20 genes, FDR $=0.003$) and CH-OH-acting oxidoreductases, and contains
22 predicted secreted proteins and 11 CAZymes; it also enriches for ``protein aggregate
centre'' (3/5). The replicate-supported nodule markers are independently dominated by
cerato-platanin family proteins and small heat-shock proteins---consistent with the
aggregate-centre enrichment recovered by a different method.

The exudophore, a newly described exudate-producing structure, yields the largest set of
replicate-supported markers of any tissue. These are consistently oxidative and secretory:
alcohol oxidase 1, glyoxal oxidase, formate and aldehyde dehydrogenases, an FAD-dependent
monooxygenase, acetate--CoA ligase, an amino-acid permease and an AA9 lytic polysaccharide
monooxygenase. Alcohol oxidase and glyoxal oxidase both generate extracellular H$_2$O$_2$, the
co-substrate required by fungal peroxidases. We report these as evidence that the corrected
pipeline recovers coherent, method-independent signal, not as a characterisation of exudophore
function; the full tissue comparison is reserved for a companion study.

\subsection{Tissue modules are not explained by ribosomal RNA load}

Because rRNA content varies among these libraries (42.0--96.6\% of assigned counts) and is
itself correlated with usable depth ($r=-0.907$ with $\log_{10}$ mRNA counts), any module
correlating with tissue could in principle reflect rRNA load rather than biology. This concern
is sharpest for the exudophore-associated \texttt{lightcyan} module, which enriches for
SSU-rRNA maturation---a ribosome-biogenesis signature that could plausibly co-vary with
ribosomal content for purely technical reasons.

We tested each module eigengene against the per-library rRNA fraction and recomputed the
tissue association as a partial correlation controlling for it
(Fig.~\ref{fig:wgcna}b, Table~\ref{tab:modules}). All six associations survive. The
\texttt{lightcyan} module correlates with rRNA fraction at only $-0.15$, so there is
essentially no technical signal to remove. Because rRNA fraction is strongly anti-correlated
with usable depth, this control also partially accounts for depth.

\section{Discussion}

\subsection{Unannotated rRNA is a silent failure mode}

The largest single component of these libraries occupied a genomic region carrying no
annotated feature, and its consequences were entirely silent. No tool raised an error.
Alignment rates looked unremarkable. The reads were reported under
\texttt{Unassigned\_NoFeatures}, a category most workflows summarise but few interrogate.

The consequence for rRNA removal deserves emphasis. Version G of the GeneLab consensus
pipeline added a parallel rRNA-removed differential expression track, implemented by dropping
features whose biotype is rRNA before renormalisation \cite{genelab7101G}. That is a sound
design, and we adopted it. But it is a \emph{feature}-based operation, and it therefore cannot
remove rRNA that was never annotated as a feature. A pipeline can execute its rRNA-removal step
faithfully, report success, and leave the majority of the ribosomal signal untouched. The
vulnerability grows precisely where it matters most---a non-model organism whose annotation was
contributed by a genome project rather than curated for expression analysis.

Two properties made this detectable: the loss was focal rather than diffuse, and the identity
was verifiable. We suggest that clustering unassigned reads and inspecting the top loci should
be routine when assignment rates are low; it costs minutes.

\subsection{The intuitive explanation was wrong, and testing it was cheap}

Low gene assignment in 3$'$-tag data has an obvious candidate explanation, and it was well
supported here. Acting on it would have been reasonable. Testing it took an afternoon and
showed it to be almost entirely wrong (Section~\ref{sec:utr}). Had we adopted the extension
without measuring its effect, we would have concluded that the annotation was now adequate,
retained a plausible-sounding correction with no benefit, and never looked for the real cause.
The generalisable point is not that UTR extension is useless---it remains correct practice for
tag data, and we retain it---but that the acceptance criterion should be specified before the
correction is applied.

\subsection{Reference choice is consequential, and standard metrics do not predict it}

The RefSeq-designated reference genome produced the worst mapping of four assemblies, by a
factor of more than four, and contiguity did not predict performance either: a contig-level
assembly outperformed a chromosome-level one on gene assignment. The property that mattered is
not reported by any standard assembly metric. For expression work on species with several
published assemblies we recommend empirical selection on assignment rate using matched
annotations, rather than selection on assembly designation, contiguity or gene count.

\subsection{What poly(A) selection did not do}\label{sec:polya}

Poly(A) selection is expected to deplete ribosomal RNA, which is not polyadenylated. Here it
did not, and the residual ribosomal fraction varied systematically with tissue. Fungal tissue
with high ribosome content presents a harder depletion problem than the cell types for which
such protocols are typically validated; where sample types of this kind are being profiled,
explicit ribosomal depletion may be the safer choice, and a pilot library is a cheap way to
establish which is needed. Because ribosomal content covaried with tissue identity, it was
confounded with the biological factor of interest, which makes the rRNA-removed analysis track
a requirement rather than a refinement---and that track could only be implemented correctly
once the missing rRNA features had been added.

\subsection{Homology-based annotation misses the interesting proteins}

Forty-one per cent of the proteome carried a Swiss-Prot match. That figure sets a hard ceiling
on every downstream enrichment analysis and explains why the two largest co-expression modules
had no enriched terms at all: they are dominated by proteins with no characterised homologue.
Lineage-specific secreted proteins are precisely the class most likely to be involved in a
novel exudate-producing structure, and precisely the class least likely to have a Swiss-Prot
homologue. Our secretome estimate is a lower bound for the same reason, and absence of
enrichment in these data is weak evidence.

\subsection{What compromised data can still support}

Roughly 2.7\% of sequenced reads reached a protein-coding feature. Establishing where the limit
falls required explicit tests rather than judgement. Including the four lowest-yield libraries
made sequencing depth the dominant axis of variation; removing them eliminated the artefact
entirely and \emph{increased} the number of genes passing expression filtering and the number
with a detectable tissue effect. Noise-dominated libraries do not merely add nothing; they
distort dispersion estimation for every other sample. Retaining all samples is not the
conservative choice.

We then applied two criteria sensitive to different failure modes---a marker test guarding
against single-replicate artefacts, and FDR-corrected module--trait association guarding
against multiplicity. They agreed completely. Where such criteria disagree, neither should be
trusted; where they converge, the conclusion is considerably stronger than either alone. Both
tissues that failed had four and three replicates respectively, more than one that passed:
sequencing depth, not replicate number, determined what was recoverable.

\subsection{The metabolic reconstruction}

The draft model is offered as a community starting point rather than a predictive tool. Its
construction exposed two limitations worth recording. Mapping enzyme commission numbers onto a
reaction database recovers core metabolism well but systematically misses steps annotated
without a complete EC, which fell disproportionately on cofactor biosynthesis; and transport
reactions, which carry no EC at all, were absent entirely from the initial draft and accounted
for most of its blocked reactions. One gapfilling case is instructive: biotin remained
unreachable, and we added it to the medium rather than to the network, because many fungi are
biotin auxotrophs and the culture medium supplies B-vitamins. Automated gapfilling will readily
invent a biosynthetic route for any metabolite declared essential; whether it should is a
biological question, not an algorithmic one.

\subsection{Limitations}

Differential expression rests on 12 libraries with two to four replicates per tissue, and only
three libraries exceeded 100,000 assigned counts. The co-expression network was constructed at
a sample size below that recommended by the method's authors \cite{langfelder2008wgcna}, and
its modules are exploratory. CAZy assignments are Pfam-derived class-level calls, not dbCAN
family assignments. The secretome is homology-transferred, not predicted \textit{de novo}.
Conclusions regarding two of the four tissues are absent not because those tissues lack
distinct biology but because these libraries could not resolve it.

\section{Conclusions}

Ribosomal RNA absent from a genome annotation is invisible to the rRNA-removal steps that
current consensus pipelines implement, and its loss is reported in a category that is rarely
examined. Reference selection for expression analysis should be empirical; the species
reference genome performed worst of four candidates here, and the property that determined
performance is not reported by standard assembly metrics. Severely compromised data can still
support conclusions, but only those that survive criteria sensitive to different failure modes,
and only after noise-dominated libraries are removed rather than retained out of caution.

\section{Methods}

\subsection{Biological material and culture}
\textit{Pleurotus ostreatus} cv.\ ``Harbor Blue P01'' was grown on micro-structured ceramic
MiniTubes (10\,cm $\times$ 5\,cm, pore size $<$300\,nm, 50\% v/v granular activated carbon)
supplied with mycoponic nutrient medium v3, at 16\,$^\circ$C, 85\% relative humidity and
CO$_2$ $\leq$1000\,ppm \cite{porterfield2026mycoponics}. Four tissue types were sampled with
four biological replicates each: exuding mycelium, fuzzy mycelium, exudophore (a newly
described exudate-producing structure), and nodule.

\noindent\textbf{[AUTHORS: sampling procedure, tissue definitions, morphological description
of the exudophore, RNA extraction method, and culture age at sampling.]}

\subsection{Sequencing}
Libraries were prepared as poly(A)-selected 3$'$-end tag libraries and sequenced on an
Illumina NovaSeq X Plus (single-end 94\,bp, 14\,nt unique molecular identifier appended to the
read name). All 16 libraries were sequenced on one flowcell and lane, so no batch term was
required. Total yield was 92,430,743 raw reads (median 5.4\,M per library).

Ribosomal RNA constituted 20--70\% of reads per library. Because this proved to be the
dominant constraint on the dataset, and because it is a property of poly(A)-selected 3$'$-tag
libraries prepared from this tissue rather than of any downstream processing choice, it is
treated here as a primary object of study rather than as a quality caveat
(Sections~\ref{sec:rrna} and~\ref{sec:polya}).

\subsection{Reference selection}
Four assemblies were compared (Table~\ref{tab:refs}); 250,000 reads per library were mapped to
genome-only indices and compared on uniquely mapped fraction, then on gene assignment rate
using matched 3$'$-extended annotations.

\subsection{Annotation construction}
Ribosomal RNA loci were identified with barrnap \cite{seemann2013barrnap} in nuclear and
mitochondrial modes and supplemented with coverage-derived blocks obtained by clustering
intergenic reads from pooled alignments. Gene 3$'$ ends were extended strand-aware by up to
500\,bp, capped at the distance to the nearest neighbouring feature minus 50\,bp. rRNA blocks
were added as \texttt{gene\_biotype "rRNA"} features on a single strand; annotating both
strands makes every overlapping read ambiguous under unstranded counting.

\subsection{Read processing and quantification}
Adapted from GL-DPPD-7101-G \cite{genelab7101G}. Reads were trimmed with fastp 1.3.6
\cite{chen2018fastp} (polyG, polyX, 3$'$ quality, minimum length 30), aligned with HISAT2 2.2.3
\cite{kim2019hisat2} (\texttt{-{}-max-intronlen 3000 -{}-rna-strandness F}), reduced to primary
alignments, deduplicated with UMI-tools 1.1.6 \cite{smith2017umitools}
(\texttt{-{}-method=directional}), and counted with featureCounts (Subread 2.1.1)
\cite{liao2014featurecounts} (\texttt{-t exon -g gene\_id -s 1}). Strandedness was established
empirically (forward: 49.6\% assigned versus 6.3\% reverse).

One deviation from the source pipeline is consequential: featureCounts replaces RSEM, because
RSEM models full-length transcript coverage and is invalid for 3$'$-tag pileups. No TPM or FPKM
values are reported anywhere, since 3$'$-tag counting yields one count per molecule independent
of transcript length.

\subsection{Differential expression}
DESeq2 1.46.0 \cite{love2014deseq2} with \verb|~ Tissue|, running two parallel tracks
(all-genes and rRNA-removed, each independently renormalised) as GL-DPPD-7101-G specifies.
Genes were filtered with \texttt{edgeR::filterByExpr} \cite{robinson2010edger}. Libraries
yielding fewer than 15,000 assigned non-rRNA counts were excluded in a sensitivity analysis.

\subsection{Functional annotation}
DIAMOND \cite{buchfink2015diamond} blastp against UniProt Swiss-Prot \cite{uniprot2023}; EC, GO
and KEGG \cite{kanehisa2000kegg} assignments derived from matched Swiss-Prot records. CAZy
classes were assigned from Pfam-A \cite{mistry2021pfam} domains detected with HMMER
\cite{eddy2011hmmer} using curated keyword rules; these are class-level, Pfam-derived calls,
not dbCAN family assignments. Secreted proteins were inferred by transferring curated SIGNAL
and TRANSMEM features from the best-hit Swiss-Prot entry.

\subsection{Co-expression network}
WGCNA 1.74 \cite{langfelder2008wgcna}, signed network, soft-threshold power 18, minimum module
size 30, merge height 0.25 (Supplementary Fig.~S1). Module--trait correlations were corrected
across all tests with Benjamini--Hochberg \cite{benjamini1995controlling}. Enrichment used the
hypergeometric test with BH correction within each module\,$\times$\,ontology.

\subsection{Metabolic reconstruction}
EC numbers were mapped onto the ModelSEED biochemistry database \cite{seaver2021modelseed} to
produce a draft with gene--protein--reaction rules traceable to supporting proteins. Transport
reactions were added and exchange bounds constrained to the culture medium. Targeted gapfilling
restricted the candidate pool by backward reachability from each blocked metabolite. Analyses
used COBRApy 0.32.1 \cite{ebrahim2013cobrapy}.

\clearpage
\section*{Figures}

\begin{figure}[htbp]\centering
\includegraphics[width=\textwidth]{fig1_read_fate.pdf}
\caption{\textbf{Read fate and ribosomal composition.}
\textbf{a}, Fate of every raw read, reconciled exactly per library.
\textbf{b}, Directly measured rRNA content, from subsampled reads mapped against the rRNA loci
themselves, independently of the counting annotation.
\textbf{c}, Assigned mRNA counts per library; dashed line marks the 15,000-count retention
threshold.}\label{fig:fate}
\end{figure}

\begin{figure}[htbp]\centering
\includegraphics[width=\textwidth]{fig2_reference_rdna.pdf}
\caption{\textbf{Reference selection and unannotated ribosomal DNA.}
\textbf{a}, Mapping outcome for four candidate assemblies; the RefSeq reference genome (PC9)
performs worst. \textbf{b}, Per-library gain in assigned counts from BOM\_ss5 relative to
PC9.15; all 16 libraries gain. \textbf{c}, Read depth across the mitochondrion. Shaded blocks
are the coverage-derived rRNA regions added here; blue marks below the axis are features
annotated in the public GTF. The highest-coverage regions carry no annotation.}\label{fig:refs}
\end{figure}

\begin{figure}[htbp]\centering
\includegraphics[width=\textwidth]{fig3_pca_retention.pdf}
\caption{\textbf{Excluding noise-dominated libraries removes a depth artefact.}
\textbf{a}, VST principal components before and after excluding four libraries below 15,000
assigned non-rRNA counts; point size is $\log_{10}$ depth. \textbf{b}, Genes passing
expression filtering and genes with a significant tissue effect both increase.}\label{fig:pca}
\end{figure}

\begin{figure}[htbp]\centering
\includegraphics[width=\textwidth]{fig4_wgcna.pdf}
\caption{\textbf{Co-expression modules.}
\textbf{a}, Module--tissue correlation; asterisks mark associations surviving
Benjamini--Hochberg correction across all 168 tests. \textbf{b}, Each surviving association
compared with its correlation to per-library rRNA fraction, and recomputed as a partial
correlation controlling for it.}\label{fig:wgcna}
\end{figure}

\clearpage
\section*{Tables}
\begin{table}[htbp]
\caption{Per-library read budget against the BOM\_ss5 reference after UMI deduplication. rRNA and mRNA are counts assigned to features of the respective biotype.}\label{tab:budget}
\footnotesize
\begin{tabular}{@{}llrrrrrc@{}}
\toprule
Tissue & Well & Raw reads & rRNA & mRNA & mRNA \% & Genes $\geq$10 &  \\
\midrule
Exuding mycelium & 1A & 8,423,109 & 701,579 & 469,981 & 5.58 & 6,017 &  \\
Exuding mycelium & 1B & 5,902,765 & 459,570 & 19,365 & 0.33 & 354 &  \\
Exuding mycelium & 1C & 2,770,639 & 64,744 & 2,620 & 0.09 & 19 & $\dagger$ \\
Exuding mycelium & 1D & 5,770,952 & 646,755 & 29,276 & 0.51 & 546 &  \\
Fuzzy mycelium & 1E & 13,369,308 & 1,032,480 & 35,940 & 0.27 & 704 &  \\
Fuzzy mycelium & 1F & 5,289,692 & 330,615 & 46,085 & 0.87 & 1,034 &  \\
Fuzzy mycelium & 1G & 6,274,968 & 714,848 & 224,213 & 3.57 & 4,507 &  \\
Fuzzy mycelium & 1H & 5,610,507 & 887,407 & 81,814 & 1.46 & 1,698 &  \\
Exudophore & 2A & 2,941,102 & 145,479 & 13,447 & 0.46 & 239 & $\dagger$ \\
Exudophore & 2B & 3,412,844 & 99,368 & 10,944 & 0.32 & 164 & $\dagger$ \\
Exudophore & 2C & 4,028,384 & 143,150 & 66,739 & 1.66 & 1,485 &  \\
Exudophore & 2D & 2,620,662 & 168,385 & 46,606 & 1.78 & 998 &  \\
Nodule & 2E & 2,431,986 & 118,075 & 7,062 & 0.29 & 109 & $\dagger$ \\
Nodule & 2F & 7,338,750 & 588,323 & 41,959 & 0.57 & 827 &  \\
Nodule & 2G & 3,473,216 & 531,658 & 40,253 & 1.16 & 813 &  \\
Nodule & 2H & 12,771,859 & 969,538 & 1,339,504 & 10.49 & 7,512 &  \\
\bottomrule
\end{tabular}
\begin{tablenotes}[flushleft]\footnotesize\item $\dagger$ excluded from the differential expression analysis (fewer than 15,000 assigned non-rRNA counts).\end{tablenotes}
\end{table}

\begin{table}[htbp]
\caption{Reference candidates compared empirically. 250{,}000 reads per library mapped to genome-only indices; values are means across all 16 libraries.}\label{tab:refs}
\footnotesize
\begin{tabular}{@{}lllrrrr@{}}
\toprule
Assembly & Accession & Level & Genes & Unique \% & Multi \% & Unaligned \% \\
\midrule
BOM\_ss5 & GCA\_056149245.1 & contig & 12,705 & 24.94 & 53.17 & 21.89 \\
BOM\_ss14 & GCA\_056149315.1 & contig & 13,310 & 24.92 & 53.18 & 21.90 \\
PC9.15 & GCA\_029852705.2 & chromosome & 13,556 & 21.49 & 53.10 & 25.42 \\
PC9 & GCF\_014466165.1 & contig & 11,849 & 5.45 & 69.14 & 25.41 \\
\bottomrule
\end{tabular}
\begin{tablenotes}[flushleft]\footnotesize\item The RefSeq-designated reference genome for the species (PC9) performed worst.\end{tablenotes}
\end{table}

\begin{table}[htbp]
\caption{Co-expression modules whose tissue association survives Benjamini--Hochberg correction across all 168 module\,$\times$\,tissue tests.}\label{tab:modules}
\footnotesize
\begin{tabular}{@{}lrlrrrrp{42mm}@{}}
\toprule
Module & Genes & Tissue & $r$ & FDR & $r_{\mathrm{rRNA}}$ & $r_{\mathrm{partial}}$ & Enriched terms \\
\midrule
greenyellow & 142 & Exudophore & -0.92 & 0.0034 & +0.28 & -0.92 & pexophagy / response to calcium ion \\
lightcyan & 99 & Exudophore & +0.91 & 0.0034 & -0.15 & +0.91 & maturation of SSU-rRNA / CENP-A containing chromatin... \\
white & 55 & Nodule & -0.88 & 0.0091 & -0.10 & -0.92 & one-carbon metabolic process \\
lightsteelblue1 & 34 & Exudophore & +0.85 & 0.0155 & -0.36 & +0.85 & chromosome segregation \\
brown & 320 & Nodule & +0.85 & 0.0155 & -0.09 & +0.85 & oxidoreductase activity / oxidoreductase activity, a... \\
saddlebrown & 49 & Exudophore & +0.84 & 0.0179 & +0.07 & +0.88 & Methane metabolism / Microbial metabolism in diverse... \\
\bottomrule
\end{tabular}
\begin{tablenotes}[flushleft]\footnotesize\item $r_{\mathrm{rRNA}}$, correlation of the module eigengene with per-library rRNA fraction; $r_{\mathrm{partial}}$, tissue correlation controlling for it.\end{tablenotes}
\end{table}


\clearpage
\section*{Supplementary Figures}
\setcounter{figure}{0}
\renewcommand{\thefigure}{S\arabic{figure}}

\begin{figure}[htbp]\centering\includegraphics[width=\textwidth]{figS1_soft_threshold.pdf}
\caption{\textbf{Co-expression network construction diagnostics.} \textbf{a}, Scale-free
topology fit against soft-threshold power; no power reaches the conventional $R^2=0.9$ at this
sample size, and 18 was used. \textbf{b}, Mean connectivity.}\end{figure}

\begin{figure}[htbp]\centering\includegraphics[width=0.62\textwidth]{figS2_utr_extension.pdf}
\caption{\textbf{Extending gene 3$'$ ends does not recover the missing reads.} Per-library
assignment rate before and after strand-aware 3$'$ extension of 91.3\% of genes by a median of
323\,bp.}\end{figure}

\begin{figure}[htbp]\centering\includegraphics[width=\textwidth]{figS3_tau_robustness.pdf}
\caption{\textbf{Specificity scores and replicate support.} \textbf{a}, Distribution of the
$\tau$ index. \textbf{b}, Of the top 50 markers per tissue, the number supported by every
replicate.}\end{figure}

\begin{figure}[htbp]\centering\includegraphics[width=0.78\textwidth]{figS4_cazy_secretome.pdf}
\caption{\textbf{Carbohydrate-active enzymes and predicted secretome.} \textbf{a}, CAZy classes
assigned from Pfam domains. \textbf{b}, Proportion of the proteome with a Swiss-Prot homologue
and, of those, the predicted secreted fraction.}\end{figure}

\begin{figure}[htbp]\centering\includegraphics[width=0.68\textwidth]{figS5_gem_stages.pdf}
\caption{\textbf{Metabolic reconstruction by stage.} Reactions able to carry flux rise from 288
in the EC-mapped draft to 2,287 after adding transport reactions, constraining exchanges to the
culture medium, and targeted cofactor gapfilling.}\end{figure}

\begin{figure}[htbp]\centering\includegraphics[width=0.78\textwidth]{figS6_duplication.pdf}
\caption{\textbf{UMI duplication rate per library.}}\end{figure}

\begin{figure}[htbp]\centering\includegraphics[width=0.82\textwidth]{figS7_intergenic_peaks.pdf}
\caption{\textbf{Unannotated loci ranked by read share.} Clustering reads that fall outside
annotated genes; the ribosomal DNA blocks (red) dominate, which is how they were located.}
\end{figure}

\clearpage
\section*{Supplementary Tables}
\begin{table}[htbp]
\caption{Sample manifest. All libraries were sequenced on one flowcell and lane.}\label{tab:s1}
\footnotesize
\begin{tabular}{@{}lll@{}}
\toprule
Library & Well & Tissue \\
\midrule
GPNJ7M\_1\_sample\_1A & 1A & Exuding mycelium \\
GPNJ7M\_2\_sample\_1B & 1B & Exuding mycelium \\
GPNJ7M\_3\_sample\_1C & 1C & Exuding mycelium \\
GPNJ7M\_4\_sample\_1D & 1D & Exuding mycelium \\
GPNJ7M\_5\_sample\_1E & 1E & Fuzzy mycelium \\
GPNJ7M\_6\_sample\_1F & 1F & Fuzzy mycelium \\
GPNJ7M\_7\_sample\_1G & 1G & Fuzzy mycelium \\
GPNJ7M\_8\_sample\_1H & 1H & Fuzzy mycelium \\
GPNJ7M\_9\_sample\_2A & 2A & Exudophore \\
GPNJ7M\_10\_sample\_2B & 2B & Exudophore \\
GPNJ7M\_11\_sample\_2C & 2C & Exudophore \\
GPNJ7M\_12\_sample\_2D & 2D & Exudophore \\
GPNJ7M\_13\_sample\_2E & 2E & Nodule \\
GPNJ7M\_14\_sample\_2F & 2F & Nodule \\
GPNJ7M\_15\_sample\_2G & 2G & Nodule \\
GPNJ7M\_16\_sample\_2H & 2H & Nodule \\
\bottomrule
\end{tabular}
\end{table}

\begin{table}[htbp]
\caption{Differential expression across all six pairwise contrasts (12 retained libraries, rRNA-removed track; $p_{\mathrm{adj}}<0.05$, $|\log_2\mathrm{FC}|>1$).}\label{tab:s2}
\footnotesize
\begin{tabular}{@{}llr@{}}
\toprule
Tissue A & Tissue B & DE genes \\
\midrule
Nodule & Fuzzy mycelium & 239 \\
Nodule & Exudophore & 198 \\
Nodule & Exuding mycelium & 74 \\
Fuzzy mycelium & Exudophore & 50 \\
Exudophore & Exuding mycelium & 36 \\
Fuzzy mycelium & Exuding mycelium & 2 \\
\bottomrule
\end{tabular}
\end{table}

\begin{table}[htbp]
\caption{Genome-scale metabolic reconstruction for \textit{P.\ ostreatus} BOM\_ss5.}\label{tab:s7}
\footnotesize
\begin{tabular}{@{}lrrrrr@{}}
\toprule
Stage & Reactions & Metabolites & Genes & Carry flux & Gapfilled \\
\midrule
EC-mapped draft & 1,915 & 2,028 & 881 & 288 & 0 \\
Transport + MNM medium & 5,243 & 5,140 & 881 & 2,262 & 0 \\
Targeted gapfilling & 5,247 & 5,140 & 881 & 2,287 & 3 \\
\bottomrule
\end{tabular}
\begin{tablenotes}[flushleft]\footnotesize\item Gapfilled reactions carry no gene association.\end{tablenotes}
\end{table}


\noindent Supplementary data files (comma/tab separated, in \texttt{latex/tables/}):
S3 replicate-supported markers; S4 co-expression module titles; S5 module enrichment;
S6 CAZyme assignments; S8 predicted secretome.
